% !TEX program = pdflatex
\documentclass[12pt]{article}
\usepackage{booktabs}
\usepackage{amsmath}
\usepackage{amssymb}
\usepackage{geometry}
\usepackage{caption}
\usepackage{threeparttable}
\usepackage{float}
\usepackage{setspace}
\usepackage{pdflscape}
\usepackage{array}
\usepackage{enumitem}
\usepackage{hyperref}

\geometry{left=1.1in, right=1.1in, top=1in, bottom=1in}
\setstretch{1.2}
\captionsetup{labelfont=bf, font=small}

\newcommand{\sym}[1]{\ensuremath{^{\text{#1}}}}
\newcommand{\se}[1]{{\small(#1)}}
\newcommand{\abs}{\textemdash}
\newcommand{\mis}{}
\newcolumntype{L}[1]{>{\raggedright\arraybackslash}p{#1}}

\title{\textbf{Outsourcing Purchase Prices}}
\date{May 2026}

\begin{document}

\maketitle

%=============================================================
\section{Regression Results}
\label{sec:regression}
%=============================================================

\subsection*{Overview}

The dependent variable in all regressions is
\(\ln p_{fjct}\), the log outsourcing purchase unit price (CNY per unit)
paid by firm \(f\) for product \(j\) in city \(c\) in year \(t\) (2017).
Three groups of explanatory variables are considered.

\begin{itemize}[leftmargin=1.5em]
  \item \textbf{Firm-level variables (\(x\)):}
    \(\ln\text{Output}_{ft}\) (log total VAT sales, size proxy);
    \(\ln K_{ft}\) (log total assets);
    \(\ln Q_{fjt}\) (log of the firm's own sales quantity of product
    \(j\)).
    \(\ln\text{Output}\) and \(\ln K\) vary at the firm\(\times\)year level
    and are absorbed by firm FE; \(\ln Q_{fjt}\) varies at the
    firm\(\times\)product level and remains identified under firm FE.

  \item \textbf{Market-level variables (\(z\)):}
    \(\ln n^{B}_{jct}\) (log number of buyers of product \(j\) in city \(c\));
    \(\ln n^{S}_{jct}\) (log number of sellers);
    \(\ln q^{m}_{jct}\) (log total market purchase quantity);
    \(\ln \bar{p}^{B}_{jct}\) (log buy-side leave-one-out market average price);
    \(\ln \bar{p}^{S}_{jct}\) (log sell-side leave-one-out market average price).
    Both leave-one-out prices subtract the firm's own buy/sell transactions
    from the city--product totals, removing the mechanical correlation between
    the dependent variable and these controls.

  \item \textbf{Product-similarity variables (\(s\)):}
    \(S_{mj}\) (input\_similarity); \(C_{mj}\) (output\_similarity). Both vary at the firm\(\times\)product level.
\end{itemize}

\paragraph{Sample.}
The sample consists of 3,410 firms randomly drawn from the VAT database.
The sample passes through the following filtering steps:
\begin{itemize}[leftmargin=1.5em,nosep]
  \item \textbf{Buy-side filter}: 3,376 of the 3,410 firms have at least one
    valid 9-digit product code in the purchase records.
  \item \textbf{Outsourcing definition}: an outsourcing product requires the
    same firm to appear on both the buy and sell side for the
    same product code. Only 2,108 firms satisfy this condition. 
  \item \textbf{Intermediary removal}: 233 firms with outsourcing ratio
    \(>90\%\) are dropped, leaving
    1,875 firms in the final regression panel.
\end{itemize}


\paragraph{Product portfolio.}
Among the 3,376 buy-side firms, each purchases on average approximately 125
distinct product types.
In the final regression panel 1,875 firms, the mean is 25.0 outsourcing products per firm
(median 3), spanning 2,143 distinct product types and 262 cities.

\paragraph{Fixed effects and identification.}
All FE specifications include firm (\(\delta_f\)), product (\(\delta_j\)), and/or city (\(\delta_c\)) fixed effects; standard errors are
clustered at the firm level.
Because the data cover a single cross-section (2017), the firm-level variables
\(\ln\text{Output}\) and \(\ln K\) are perfectly collinear with firm FE and are
identified only in OLS specifications. 

\bigskip

%-------------------------------------------------------------
\subsection{T1: Firm-Level Variables}
\label{sec:t1}
%-------------------------------------------------------------

The baseline equation is
\begin{equation}
  \ln p_{fjct}
  = \gamma_1 \ln \text{Output}_{ft}
  + \gamma_2 \ln K_{ft}
  + \gamma_3 \ln Q_{fjt}
  + \delta_f + \delta_j + \delta_c
  + \varepsilon_{fjct}.
  \label{eq:t1}
\end{equation}

\begin{table}[H]
  \centering
  \caption{Variable Definitions --- Table~T1}
  \label{tab:vardef_t1}
  \begin{tabular}{L{3.5cm} L{10.5cm}}
    \toprule
    Variable & Definition \\
    \midrule
    \(p_{fjct}\) &
       Log outsourcing purchase unit price,
      firm \(f\), product \(j\), city \(c\), year \(t\). \\[4pt]
    \(\ln \text{Output}_{ft}\) &
      Log annual VAT sales\\[4pt]
    \(\ln K_{ft}\) &
      Log total assets
     \\[4pt]
    \(\ln Q_{fjt}\) &
      Log of the firm's own output quantity of product \(j\)\\[4pt]
    \(\delta_f,\,\delta_j,\,\delta_c\) &
      Firm, product, and city fixed effects. \\
    \bottomrule
  \end{tabular}
\end{table}

\paragraph{Key findings.}
\begin{itemize}[leftmargin=1.5em,nosep]
  \item In OLS, firm size \(\ln\text{Output}\) is significantly positive
    (\(+0.234^{***}\)): larger firms pay higher outsourcing prices (maybe better quality?); capital
    \(\ln K\) is significantly negative (\(-0.156^{***}\)): more capital-rich
    firms pay lower prices.
  \item \(\ln Q\) is significantly negative in both OLS and the firm-FE column
    (\(-0.177^{***}\), \(-0.169^{***}\)): larger production quantities go with
    lower outsourcing prices, consistent with quantity discounts.
\end{itemize}

\begin{table}[H]
  \centering
  \caption{T1: Firm-Level Variables Only}
  \label{tab:t1}
  \begin{threeparttable}
    \begin{tabular}{lcc}
      \toprule
      & (1) OLS & (2) +Firm FE \\
      \midrule
      \(\ln \text{Output}\) & \(0.234\sym{***}\) & \abs \\
                             & \se{0.051}         & \\[4pt]
      \(\ln K\)              & \(-0.156\sym{***}\) & \abs \\
                             & \se{0.047}          & \\[4pt]
      \(\ln Q\)              & \(-0.177\sym{***}\) & \(-0.169\sym{***}\) \\
                             & \se{0.013}          & \se{0.013} \\[4pt]
      Cons               & \(3.527\sym{***}\) & \(4.969\sym{***}\) \\
                             & \se{0.452}         & \se{0.081} \\
      \midrule
      Firm FE     & No  & Yes \\
      Product FE  & No  & No  \\
      City FE     & No  & No  \\
      \(N\)       & 29,247 & 28,915 \\
      Adj.\ \(R^2\) & 0.081 & 0.311 \\
      \bottomrule
    \end{tabular}
  \end{threeparttable}
\end{table}

%-------------------------------------------------------------
\subsection{T2: Firm-Level Variables and Product Similarity}
\label{sec:t2}
%-------------------------------------------------------------

The equation adds product-similarity terms:
\begin{equation}
  \ln p_{fjct}
  = \gamma_1 \ln \text{Output}_{ft}
  + \gamma_2 \ln K_{ft}
  + \gamma_3 \ln Q_{fjt}
  + \beta_1 S_{mj}
  + \beta_2 C_{mj}
  + \delta_f + \delta_j + \delta_c
  + \varepsilon_{fjct}.
  \label{eq:t2}
\end{equation}

\begin{table}[H]
  \centering
  \caption{Variable Definitions --- Table~T2 }
  \label{tab:vardef_t2}
  \begin{tabular}{L{3.5cm} L{10.5cm}}
    \toprule
    Variable & Definition \\
    \midrule
    \(S_{mj}\) &
      Input similarity. \\[4pt]
    \(C_{mj}\) &
      Output similarity \\[4pt]
    \(m\) &
      Core product: with the highest net production. \\
    \bottomrule
  \end{tabular}
\end{table}

\paragraph{Key findings.}
\begin{itemize}[leftmargin=1.5em,nosep]
  \item In OLS, firm size is positive (\(+0.267^{***}\)) and capital is negative
    (\(-0.112^{***}\)), as in T1.
  \item Similarity \(S_{mj}\), \(C_{mj}\) are strongly positive in OLS
    (\(+1.491^{***}\), \(+1.481^{***}\)) and attenuate sharply under Product FE
    to marginal significance (\(0.186^{*}\), \(0.263^{*}\)): the similarity
    effect is largely explained by product-level heterogeneity.
  \item \(\ln Q\) is significantly negative from OLS (\(-0.226^{***}\)) to full
    FE (\(-0.070^{***}\)).
\end{itemize}

\begin{table}[H]
  \centering
  \caption{T2: Firm-Level Variables and Product Similarity}
  \label{tab:t2}
  \begin{threeparttable}
    \begin{tabular}{lcccc}
      \toprule
      & (1) OLS & (2) +Firm FE & (3) +Firm+Prod & (4) +Firm+Prod+City \\
      \midrule
      \(\ln \text{Output}\) & \(0.267\sym{***}\) & \abs & \abs & \abs \\
                             & \se{0.042} & & & \\[4pt]
      \(\ln K\) & \(-0.112\sym{***}\) & \abs & \abs & \abs \\
                & \se{0.040} & & & \\[4pt]
      \(\ln Q\) & \(-0.226\sym{***}\) & \(-0.193\sym{***}\) & \(-0.070\sym{***}\) & \(-0.070\sym{***}\) \\
                & \se{0.013} & \se{0.014} & \se{0.008} & \se{0.008} \\[4pt]
      \(S_{mj}\) & \(1.491\sym{***}\) & \(0.460\sym{**}\) & \(0.186\sym{*}\) & \(0.186\sym{*}\) \\
                 & \se{0.255} & \se{0.186} & \se{0.109} & \se{0.109} \\[4pt]
      \(C_{mj}\) & \(1.481\sym{***}\) & \(1.361\sym{***}\) & \(0.263\sym{*}\) & \(0.263\sym{*}\) \\
                 & \se{0.285} & \se{0.210} & \se{0.141} & \se{0.141} \\[4pt]
      Cons & \(1.826\sym{***}\) & \(4.846\sym{***}\) & \(4.274\sym{***}\) & \(4.274\sym{***}\) \\
               & \se{0.452} & \se{0.076} & \se{0.047} & \se{0.048} \\
      \midrule
      Firm FE    & No  & Yes & Yes & Yes \\
      Product FE & No  & No  & Yes & Yes \\
      City FE    & No  & No  & No  & Yes \\
      \(N\)      & 29,237 & 28,905 & 28,552 & 28,552 \\
      Adj.\ \(R^2\) & 0.127 & 0.321 & 0.595 & 0.592 \\
      \bottomrule
    \end{tabular}
  \end{threeparttable}
\end{table}

%-------------------------------------------------------------
\subsection{T3: Market-Level Variables Only}
\label{sec:t3}
%-------------------------------------------------------------

The market-level specification is
\begin{equation}
  \ln p_{fjct}
  = \lambda_1 \ln n^{B}_{jct}
  + \lambda_2 \ln n^{S}_{jct}
  + \lambda_3 \ln q^{m}_{jct}
  + \lambda_4 \ln \bar{p}^{B}_{jct}
  + \lambda_5 \ln \bar{p}^{S}_{jct}
  + \delta_f + \delta_j + \delta_c
  + \varepsilon_{fjct}.
  \label{eq:t3}
\end{equation}

\begin{table}[H]
  \centering
  \caption{Variable Definitions --- Table~T3}
  \label{tab:vardef_t3}
  \begin{tabular}{L{3.5cm} L{10.5cm}}
    \toprule
    Variable & Definition \\
    \midrule
    \(\ln n^{B}_{jct}\) &
      Log number of buyers of product \(j\) in city \(c\)  \\[4pt]
    \(\ln n^{S}_{jct}\) &
      Log number of sellers of product \(j\) in city \(c\) \\[4pt]
    \(\ln q^{m}_{jct}\) &
      Log total market purchase quantity of product \(j\) in city \(c\) \\[4pt]
    \(\ln \bar{p}^{B}_{jct}\) &
      Log buy-side leave-one-out market average price  \\[4pt]
    \(\ln \bar{p}^{S}_{jct}\) &
      Log sell-side leave-one-out market average price  \\
    \bottomrule
  \end{tabular}
\end{table}

\paragraph{Key findings.}
\begin{itemize}[leftmargin=1.5em,nosep]
  \item The contrast between the two leave-one-out prices is the key new
    finding. The buy-side price \(\ln\bar{p}^{B}\) is strongly positive in
    OLS/Firm FE (\(0.440^{***}\), \(0.365^{***}\)) but collapses to
    \(-0.029^{*}\) under Product FE---the earlier strong correlation was driven
    by product composition and mechanical correlation.
  \item The sell-side price \(\ln\bar{p}^{S}\), by contrast, remains positive and
    significant even under full FE (\(0.052^{***}\)), carrying price information
    beyond product FE.
  \item The buyer count \(\ln n^{B}\) is robustly positive (\(0.115^{***}\)), and
    market volume \(\ln q^{m}\) turns negative under Product FE
    (\(-0.100^{***}\)).
\end{itemize}


\begin{table}[H]
  \centering
  \caption{T3: Market-Level Variables Only}
  \label{tab:t3}
  \begin{threeparttable}
    \begin{tabular}{lcccc}
      \toprule
      & (1) OLS & (2) +Firm FE & (3) +Firm+Prod & (4) +Firm+Prod+City \\
      \midrule
      \(\ln n^{B}\) & \(0.170\sym{***}\) & \(0.098\sym{***}\) & \(0.115\sym{***}\) & \(0.115\sym{***}\) \\
                    & \se{0.051} & \se{0.029} & \se{0.040} & \se{0.040} \\[4pt]
      \(\ln n^{S}\) & \(0.104\sym{*}\) & \(0.064\sym{*}\) & \(-0.006\) & \(-0.006\) \\
                    & \se{0.061} & \se{0.036} & \se{0.035} & \se{0.035} \\[4pt]
      \(\ln q^{m}\) & \(-0.009\) & \(-0.029\sym{*}\) & \(-0.100\sym{***}\) & \(-0.100\sym{***}\) \\
                    & \se{0.017} & \se{0.015} & \se{0.017} & \se{0.017} \\[4pt]
      \(\ln \bar{p}^{B}\) & \(0.440\sym{***}\) & \(0.365\sym{***}\) & \(-0.029\sym{*}\) & \(-0.029\sym{*}\) \\
                          & \se{0.017} & \se{0.016} & \se{0.017} & \se{0.017} \\[4pt]
      \(\ln \bar{p}^{S}\) & \(0.302\sym{***}\) & \(0.246\sym{***}\) & \(0.052\sym{***}\) & \(0.052\sym{***}\) \\
                          & \se{0.020} & \se{0.016} & \se{0.011} & \se{0.011} \\[4pt]
      Cons & \(0.218\) & \(1.643\sym{***}\) & \(4.912\sym{***}\) & \(4.912\sym{***}\) \\
               & \se{0.259} & \se{0.208} & \se{0.302} & \se{0.302} \\
      \midrule
      Firm FE    & No  & Yes & Yes & Yes \\
      Product FE & No  & No  & Yes & Yes \\
      City FE    & No  & No  & No  & Yes \\
      \(N\)      & 46,444 & 45,956 & 45,669 & 45,669 \\
      Adj.\ \(R^2\) & 0.310 & 0.468 & 0.604 & 0.602 \\
      \bottomrule
    \end{tabular}
  \end{threeparttable}
\end{table}

%-------------------------------------------------------------
\subsection{T4: Firm-Level and Market-Level Variables}
\label{sec:t4}
%-------------------------------------------------------------

The combined firm and market specification is
\begin{equation}
  \ln p_{fjct}
  = \boldsymbol{\gamma}' \mathbf{x}_{ft}
  + \boldsymbol{\lambda}' \mathbf{z}_{jct}
  + \delta_f + \delta_j + \delta_c
  + \varepsilon_{fjct},
  \label{eq:t4}
\end{equation}
where \(\mathbf{x}\) includes the production quantity \(\ln Q_{fjt}\).

\paragraph{Key findings.}
\begin{itemize}[leftmargin=1.5em,nosep]
  \item In OLS, firm size is positive (\(+0.157^{***}\)) and capital is negative
    (\(-0.101^{***}\)); \(\ln Q\) is significantly negative from OLS
    (\(-0.126^{***}\)) to full FE (\(-0.061^{***}\)).
  \item The buy-side price \(\ln\bar{p}^{B}\) is significant in OLS/Firm FE
    (\(0.503^{***}\), \(0.403^{***}\)) but drops to \(-0.008\) under Product FE,
    whereas the sell-side price \(\ln\bar{p}^{S}\) stays positive and significant
    under full FE (\(0.047^{***}\)).
  \item Market volume \(\ln q^{m}\) is positive in OLS (\(0.077^{***}\)) but
    negative under full FE (\(-0.078^{***}\)); the buyer count is positive under
    full FE (\(0.106^{**}\)).
\end{itemize}

\begin{table}[H]
  \centering
  \caption{T4: Firm-Level and Market-Level Variables}
  \label{tab:t4}
  \begin{threeparttable}
    \begin{tabular}{lcccc}
      \toprule
      & (1) OLS & (2) +Firm FE & (3) +Firm+Prod & (4) +Firm+Prod+City \\
      \midrule
      \(\ln \text{Output}\) & \(0.157\sym{***}\) & \abs & \abs & \abs \\
                             & \se{0.037} & & & \\[4pt]
      \(\ln K\) & \(-0.101\sym{***}\) & \abs & \abs & \abs \\
                & \se{0.033} & & & \\[4pt]
      \(\ln Q\) & \(-0.126\sym{***}\) & \(-0.122\sym{***}\) & \(-0.061\sym{***}\) & \(-0.061\sym{***}\) \\
                & \se{0.010} & \se{0.009} & \se{0.008} & \se{0.008} \\[4pt]
      \(\ln n^{B}\) & \(0.045\) & \(0.132\sym{***}\) & \(0.106\sym{**}\) & \(0.106\sym{**}\) \\
                    & \se{0.060} & \se{0.032} & \se{0.052} & \se{0.052} \\[4pt]
      \(\ln n^{S}\) & \(0.138\sym{*}\) & \(-0.008\) & \(0.012\) & \(0.012\) \\
                    & \se{0.075} & \se{0.040} & \se{0.040} & \se{0.040} \\[4pt]
      \(\ln q^{m}\) & \(0.077\sym{***}\) & \(0.033\sym{**}\) & \(-0.078\sym{***}\) & \(-0.078\sym{***}\) \\
                    & \se{0.021} & \se{0.016} & \se{0.022} & \se{0.022} \\[4pt]
      \(\ln \bar{p}^{B}\) & \(0.503\sym{***}\) & \(0.403\sym{***}\) & \(-0.008\) & \(-0.008\) \\
                          & \se{0.023} & \se{0.019} & \se{0.021} & \se{0.021} \\[4pt]
      \(\ln \bar{p}^{S}\) & \(0.242\sym{***}\) & \(0.219\sym{***}\) & \(0.047\sym{***}\) & \(0.047\sym{***}\) \\
                          & \se{0.022} & \se{0.019} & \se{0.015} & \se{0.015} \\[4pt]
      Cons & \(-0.938\sym{*}\) & \(1.385\sym{***}\) & \(4.664\sym{***}\) & \(4.664\sym{***}\) \\
               & \se{0.503} & \se{0.241} & \se{0.363} & \se{0.364} \\
      \midrule
      Firm FE    & No  & Yes & Yes & Yes \\
      Product FE & No  & No  & Yes & Yes \\
      City FE    & No  & No  & No  & Yes \\
      \(N\)      & 28,967 & 28,636 & 28,287 & 28,287 \\
      Adj.\ \(R^2\) & 0.332 & 0.472 & 0.597 & 0.595 \\
      \bottomrule
    \end{tabular}
  \end{threeparttable}
\end{table}

%-------------------------------------------------------------
\subsection{T5: Full Specification}
\label{sec:t5}
%-------------------------------------------------------------

The full specification combines all three variable groups:
\begin{equation}
  \ln p_{fjct}
  = \boldsymbol{\gamma}' \mathbf{x}_{ft}
  + \boldsymbol{\beta}' \mathbf{s}_{fj}
  + \boldsymbol{\lambda}' \mathbf{z}_{jct}
  + \delta_f + \delta_j + \delta_c
  + \varepsilon_{fjct}.
  \label{eq:t5}
\end{equation}

\paragraph{Key findings.}
\begin{itemize}[leftmargin=1.5em,nosep]
  \item In OLS, firm size is positive (\(+0.181^{***}\)) and capital is negative
    (\(-0.074^{**}\)); similarity \(S_{mj}\), \(C_{mj}\) are strongly positive in
    OLS (\(1.076^{***}\), \(0.888^{***}\)) and attenuate to marginal
    significance under Product FE (\(0.197^{*}\), \(0.258^{*}\)).
  \item The buy-side price drops to zero under full FE (\(-0.011\)) while the
    sell-side price stays positive (\(\ln\bar{p}^{S}=0.047^{***}\)).
  \item Buyer count (\(0.102^{**}\)), market volume (\(-0.082^{***}\)), and
    production quantity (\(-0.068^{***}\)) all retain explanatory content under
    full FE.
\end{itemize}

\begin{table}[H]
  \centering
  \caption{T5: Full Specification}
  \label{tab:t5}
  \begin{threeparttable}
    \begin{tabular}{lcccc}
      \toprule
      & (1) OLS & (2) +Firm FE & (3) +Firm+Prod & (4) +Firm+Prod+City \\
      \midrule
      \(\ln \text{Output}\) & \(0.181\sym{***}\) & \abs & \abs & \abs \\
                             & \se{0.033} & & & \\[4pt]
      \(\ln K\) & \(-0.074\sym{**}\) & \abs & \abs & \abs \\
                & \se{0.029} & & & \\[4pt]
      \(\ln Q\) & \(-0.158\sym{***}\) & \(-0.138\sym{***}\) & \(-0.068\sym{***}\) & \(-0.068\sym{***}\) \\
                & \se{0.010} & \se{0.009} & \se{0.008} & \se{0.008} \\[4pt]
      \(S_{mj}\) & \(1.076\sym{***}\) & \(0.386\sym{***}\) & \(0.197\sym{*}\) & \(0.197\sym{*}\) \\
                 & \se{0.185} & \se{0.115} & \se{0.107} & \se{0.108} \\[4pt]
      \(C_{mj}\) & \(0.888\sym{***}\) & \(0.795\sym{***}\) & \(0.258\sym{*}\) & \(0.258\sym{*}\) \\
                 & \se{0.229} & \se{0.149} & \se{0.139} & \se{0.140} \\[4pt]
      \(\ln n^{B}\) & \(0.092\) & \(0.135\sym{***}\) & \(0.102\sym{**}\) & \(0.102\sym{**}\) \\
                    & \se{0.059} & \se{0.031} & \se{0.052} & \se{0.052} \\[4pt]
      \(\ln n^{S}\) & \(0.107\) & \(-0.003\) & \(0.015\) & \(0.015\) \\
                    & \se{0.074} & \se{0.041} & \se{0.040} & \se{0.040} \\[4pt]
      \(\ln q^{m}\) & \(0.048\sym{**}\) & \(0.025\) & \(-0.082\sym{***}\) & \(-0.082\sym{***}\) \\
                    & \se{0.020} & \se{0.016} & \se{0.022} & \se{0.022} \\[4pt]
      \(\ln \bar{p}^{B}\) & \(0.467\sym{***}\) & \(0.393\sym{***}\) & \(-0.011\) & \(-0.011\) \\
                          & \se{0.022} & \se{0.019} & \se{0.021} & \se{0.021} \\[4pt]
      \(\ln \bar{p}^{S}\) & \(0.228\sym{***}\) & \(0.215\sym{***}\) & \(0.047\sym{***}\) & \(0.047\sym{***}\) \\
                          & \se{0.021} & \se{0.019} & \se{0.015} & \se{0.015} \\[4pt]
      Cons & \(-1.602\sym{***}\) & \(1.413\sym{***}\) & \(4.712\sym{***}\) & \(4.712\sym{***}\) \\
               & \se{0.491} & \se{0.236} & \se{0.363} & \se{0.364} \\
      \midrule
      Firm FE    & No  & Yes & Yes & Yes \\
      Product FE & No  & No  & Yes & Yes \\
      City FE    & No  & No  & No  & Yes \\
      \(N\)      & 28,957 & 28,626 & 28,277 & 28,277 \\
      Adj.\ \(R^2\) & 0.352 & 0.476 & 0.598 & 0.595 \\
      \bottomrule
    \end{tabular}
  \end{threeparttable}
\end{table}

%-------------------------------------------------------------
\subsection{T6: Interactions --- Firm Size \texorpdfstring{$\times$}{x} Market Conditions}
\label{sec:t6}
%-------------------------------------------------------------

To explore heterogeneous price-setting, all specifications include firm, product,
and city FE; \(\ln\text{Output}\) enters only through its interactions with the
market variables. The production quantity \(\ln Q_{fjt}\) is always included.
The equation is
\begin{equation}
  \ln p_{fjct}
  = \gamma_3 \ln Q_{fjt}
  + \boldsymbol{\lambda}' \mathbf{z}_{jct}
  + \sum_{k} \mu_k \bigl(\ln \text{Output}_{ft} \times z^{k}_{jct}\bigr)
  + \delta_f + \delta_j + \delta_c
  + \varepsilon_{fjct}.
  \label{eq:t6}
\end{equation}
The interaction terms vary at the firm\(\times\)product level and remain
identified under all three FE.

\begin{table}[H]
  \centering
  \caption{Variable Definitions --- Table~T6 (interaction terms)}
  \label{tab:vardef_t6}
  \begin{tabular}{L{3.8cm} L{10.2cm}}
    \toprule
    Variable & Definition \\
    \midrule
    \(\ln\text{Output} \times \ln n^{B}\) &
      Size \(\times\) buyer-market thickness.  \\[4pt]
    \(\ln\text{Output} \times \ln n^{S}\) &
      Size \(\times\) seller competition.  \\[4pt]
    \(\ln\text{Output} \times \ln\bar{p}^{B}\) &
      Size \(\times\) buy-side market price level. \\[4pt]
    \(\ln\text{Output} \times \ln\bar{p}^{S}\) &
      Size \(\times\) sell-side market price level. \\[4pt]
    \(\ln\text{Output} \times \ln q^{m}\) &
      Size \(\times\) market volume.  \\
    \bottomrule
  \end{tabular}
\end{table}

\paragraph{Key findings.}
\begin{itemize}[leftmargin=1.5em,nosep]
  \item All columns are full-FE specifications; \(\ln\text{Output}\) enters only
    through the interaction terms, so its main effect is not separately
    estimated.
  \item The single-interaction columns reveal size heterogeneity: larger firms
    pay lower prices in markets with more buyers (\(\hat\mu_1 = -0.019^{***}\))
    and more sellers (\(\hat\mu_2 = -0.023^{***}\)), and are more sensitive to
    both the buy-side and sell-side market prices (\(+0.011^{***}\) each; the
    corresponding main effects turn negative when their interaction is added).
  \item In the full specification (column 6), multicollinearity renders most
    single interactions insignificant; only the size--volume interaction
    (\(-0.011^{***}\)) remains robust. \(\ln Q\) is robustly negative across all
    columns.
\end{itemize}

\begin{table}[H]
  \centering
  \caption{T6: Interactions --- Firm Size \(\times\) Market Conditions}
  \label{tab:t6}
  \begin{threeparttable}
    \begin{tabular}{lcccccc}
      \toprule
      & (1) Baseline & (2) +Size\(\times\)nB & (3) +Size\(\times\)nS
        & (4) +Size\(\times\)\(\bar{p}^B\) & (5) +Size\(\times\)\(\bar{p}^S\) & (6) All \\
      \midrule
      \(\ln Q\) & \(-0.047\sym{***}\) & \(-0.047\sym{***}\) & \(-0.047\sym{***}\) & \(-0.047\sym{***}\) & \(-0.047\sym{***}\) & \(-0.046\sym{***}\) \\
                & \se{0.006} & \se{0.006} & \se{0.006} & \se{0.006} & \se{0.006} & \se{0.006} \\[4pt]
      \(\ln n^{B}\) & \(0.133\sym{***}\) & \(0.516\sym{***}\) & \(0.138\sym{***}\) & \(0.138\sym{***}\) & \(0.138\sym{***}\) & \(0.009\) \\
                    & \se{0.040} & \se{0.101} & \se{0.040} & \se{0.040} & \se{0.040} & \se{0.148} \\[4pt]
      \(\ln n^{S}\) & \(-0.003\) & \(-0.010\) & \(0.457\sym{***}\) & \(-0.005\) & \(-0.006\) & \(0.211\) \\
                    & \se{0.035} & \se{0.035} & \se{0.115} & \se{0.034} & \se{0.035} & \se{0.156} \\[4pt]
      \(\ln q^{m}\) & \(-0.083\sym{***}\) & \(-0.084\sym{***}\) & \(-0.085\sym{***}\) & \(-0.087\sym{***}\) & \(-0.085\sym{***}\) & \(0.134\sym{*}\) \\
                    & \se{0.017} & \se{0.017} & \se{0.017} & \se{0.017} & \se{0.017} & \se{0.072} \\[4pt]
      \(\ln \bar{p}^{B}\) & \(-0.018\) & \(-0.018\) & \(-0.019\) & \(-0.236\sym{***}\) & \(-0.017\) & \(0.083\) \\
                          & \se{0.017} & \se{0.017} & \se{0.017} & \se{0.073} & \se{0.017} & \se{0.087} \\[4pt]
      \(\ln \bar{p}^{S}\) & \(0.052\sym{***}\) & \(0.052\sym{***}\) & \(0.052\sym{***}\) & \(0.052\sym{***}\) & \(-0.161\sym{**}\) & \(-0.075\) \\
                          & \se{0.011} & \se{0.011} & \se{0.011} & \se{0.011} & \se{0.064} & \se{0.074} \\[4pt]
      Size\(\times\)nB & \mis & \(-0.019\sym{***}\) & \mis & \mis & \mis & \(0.007\) \\
                       &     & \se{0.005}          &     &     &     & \se{0.007} \\[4pt]
      Size\(\times\)nS & \mis & \mis & \(-0.023\sym{***}\) & \mis & \mis & \(-0.011\) \\
                       &     &     & \se{0.006}           &     &     & \se{0.008} \\[4pt]
      Size\(\times\)\(\bar{p}^B\) & \mis & \mis & \mis & \(0.011\sym{***}\) & \mis & \(-0.005\) \\
                                  &     &     &     & \se{0.004}          &     & \se{0.004} \\[4pt]
      Size\(\times\)\(\bar{p}^S\) & \mis & \mis & \mis & \mis & \(0.011\sym{***}\) & \(0.006\) \\
                                  &     &     &     &     & \se{0.003}          & \se{0.004} \\[4pt]
      Size\(\times\)qm & \mis & \mis & \mis & \mis & \mis & \(-0.011\sym{***}\) \\
                       &     &     &     &     &     & \se{0.004} \\[4pt]
      Cons & \(4.749\sym{***}\) & \(4.702\sym{***}\) & \(4.726\sym{***}\) & \(4.786\sym{***}\) & \(4.773\sym{***}\) & \(4.744\sym{***}\) \\
               & \se{0.298} & \se{0.291} & \se{0.290} & \se{0.299} & \se{0.299} & \se{0.290} \\
      \midrule
      Firm FE    & Yes & Yes & Yes & Yes & Yes & Yes \\
      Product FE & Yes & Yes & Yes & Yes & Yes & Yes \\
      City FE    & Yes & Yes & Yes & Yes & Yes & Yes \\
      \(N\)      & 45,669 & 45,669 & 45,669 & 45,669 & 45,669 & 45,669 \\
      Adj.\ \(R^2\) & 0.605 & 0.605 & 0.605 & 0.605 & 0.605 & 0.606 \\
      \bottomrule
    \end{tabular}
  \end{threeparttable}
\end{table}

%=============================================================
\section{Results versus Model Predictions}
%=============================================================

This section confronts the estimates with the model's predictions, separating the
results that match expectations from those that diverge, and flags the
identification issue most worth addressing.

\subsection*{In line with predictions}

\begin{itemize}[leftmargin=1.5em,nosep]
  \item \textbf{Quantity discount.} The production quantity \(\ln Q\) is robustly
    negative across all specifications (OLS \(-0.177^{***}\), full FE
    \(-0.047^{***}\)), consistent with the expectation that larger purchase
    volumes command lower unit prices.
  \item \textbf{Size \(\times\) market-thickness interactions are negative.}
    Larger firms pay lower prices in markets with more buyers
    (\(\mu_1=-0.019^{***}\)), more sellers (\(\mu_2=-0.023^{***}\)), and larger
    total volume (Size\(\times q^m=-0.011^{***}\)), matching the prediction that
    size confers bargaining power.
  \item \textbf{Sell-side price carries independent information.} The sell-side
    leave-one-out price \(\ln\bar{p}^{S}\) stays significantly positive under full
    FE (\(0.052^{***}\)), consistent with same-side market price levels passing
    through to a firm's purchase price.
\end{itemize}

\subsection*{Diverging from predictions}

\begin{itemize}[leftmargin=1.5em,nosep]
  \item \textbf{Firm-size main effect has the wrong sign (the key divergence).}
    The model predicts that larger firms, with stronger bargaining power, pay
    lower prices; but in OLS \(\ln\text{Output}\) is significantly positive
    (\(+0.234^{***}\)). Crucially, the size\(\times\)market interactions are
    \emph{negative}, meaning within a given market larger firms do bargain prices
    down --- only the size \emph{main effect} (across-market comparison) is
    positive. The positive sign therefore more likely reflects a
    quality/positioning selection effect (larger firms source higher-quality
    inputs) rather than bargaining power, so
    \(\ln\text{Output}\) should not be read as a proxy for bargaining power.
  \item \textbf{Product similarity has the wrong sign.} Theory predicts that the
    more similar the outsourced product is to the firm's own (technological
    proximity, less information asymmetry, even the option to self-produce), the
    lower the purchase price; yet \(S_{mj}\) and \(C_{mj}\) are strongly positive
    in OLS (\(+1.491^{***}\), \(+1.481^{***}\)) and decay to barely significant
    once Product FE are added (\(0.186^{*}\), \(0.263^{*}\)). 
  \item \textbf{Seller count is insignificant.} The model predicts that more
    sellers (greater supply-side competition) lower prices, yet \(\ln n^S\) is
    insignificant in the main regressions, showing no such competition effect.
  \item \textbf{Buyer count is positive.} \(\ln n^B\) is robustly positive across
    specifications (\(0.115^{***}\)), opposite to the naive prediction that more
    demand weakens bargaining; it would need reinterpreting as a demand-congestion
    or market-buoyancy effect.
  \item \textbf{Market volume is sign-unstable.} \(\ln q^m\) is positive in OLS
    (\(+0.077^{***}\)) but turns significantly negative once Product FE are added
    (\(-0.100^{***}\)); the sign flips with the specification, indicating that
    what it captures differs before and after controlling for product mix.
\end{itemize}

\paragraph{Data caveat.}
One important caveat: the price \(p_{fjct}\) is computed as invoice value divided
by quantity, but the \emph{units of quantity are not yet harmonized} (different
firms and products may report in kilograms, tonnes, pieces, boxes, etc.).
Inconsistent units make unit prices within the same 9-digit product not directly
comparable and inject noise into the coefficients above (especially the sign and
magnitude of the level effects). A follow-up should harmonize the units---or
recompute unit prices at a finer product\(\times\)unit level---and re-check, to
confirm that the divergences in this section are driven by economic mechanisms
rather than by a units-of-measure artifact.

%=============================================================
\section{SQL Data Extraction Logic}
\label{sec:sql}
%=============================================================

\subsection*{Background}

Raw data originate from the national VAT invoice database for 2017, stored in
a relational database. Five aggregated views were extracted and saved as flat
files for downstream Python processing. Table~\ref{tab:sql_views} summarizes
each view.

\begin{table}[H]
  \centering
  \caption{Extracted Data Views from the VAT Database}
  \label{tab:sql_views}
  \begin{threeparttable}
    \begin{tabular}{L{2.8cm} L{4.2cm} L{3.0cm} L{4.2cm}}
      \toprule
      View name & Unit of observation & Raw / Clean rows & Contents \\
      \midrule
      \texttt{firm\_buy} &
        Buyer firm \(\times\) product code &
        473,268 / 425,713 &
        Purchase amount and quantity; 2,731 distinct products. \\[4pt]
      \texttt{firm\_sell} &
        Seller firm \(\times\) product code &
        101,844 / 89,584 &
        Sales amount and quantity; 2,542 distinct products. Sales quantity
        feeds \(\ln Q_{fjt}\) and the sell-side leave-one-out price. \\[4pt]
      \texttt{firm\_city} &
        Firm ID \(\times\) prefecture city &
        3,410 / 3,410 &
        City assignment for each sample firm. Used to merge city
        into purchase/sale files which contain no city field. \\[4pt]
      \texttt{city\_buy} &
        Buyer prefecture \(\times\) product code &
        2,104,901 / 1,393,329 &
        Universe-level buyer count, amount, and quantity. Used to
        construct \(n^{B}\), \(q^{m}\), and the leave-one-out
        \(\bar{p}^{B}\). \\[4pt]
      \texttt{city\_sell} &
        Seller prefecture \(\times\) product code &
        1,540,065 / 954,222 &
        Universe-level seller count, amount, and quantity. Used to construct
        \(n^{S}\) and the leave-one-out \(\bar{p}^{S}\). \\
      \bottomrule
    \end{tabular}
  \end{threeparttable}
\end{table}

\subsection*{Core Python Processing Steps}

\begin{enumerate}[leftmargin=1.8em, label=\arabic*.]

  \item \textbf{Product code standardization.}
    Raw codes are right-padded to 19 digits; the first 9 digits are retained as
    the standardized product code. Only the 2,778 valid codes from the national
    product catalog are kept.

  \item \textbf{Outsourcing identification.}
    A firm-product pair is classified as an outsourcing transaction if and only
    if the firm appears in \emph{both} \texttt{firm\_buy} and
    \texttt{firm\_sell} for the same product code (i.e., the firm both purchases
    and sells the same product), with
    \(\min(\text{purchase value},\, \text{sales value}) > 0\).

  \item \textbf{City merge.}
    Because the purchase and sale files contain no geographic identifier,
    \texttt{firm\_city} is joined to both files on firm ID to assign each
    transaction to a prefecture.

  \item \textbf{Market condition construction.}
    Universe views \texttt{city\_buy} and \texttt{city\_sell} are aggregated at
    the product\(\times\)city level to produce \(n^{B}_{jct}\),
    \(n^{S}_{jct}\), \(q^{m}_{jct}\), \(\bar{p}^{B}_{jct}\), and
    \(\bar{p}^{S}_{jct}\). Both market average prices are constructed on a
    leave-one-out basis: each firm's own value and quantity (on the buy and sell
    side respectively) are subtracted from the city--product totals before
    taking the ratio, so the controls do not mechanically contain the firm's own
    price (firms that are the sole buyer/seller are set to missing).
    These aggregates are \emph{not} restricted to sample firms, ensuring that
    the market measures capture the full competitive environment.

  \item \textbf{Intermediary removal.}
    Firms flagged as intermediaries (\texttt{is\_intermediary == 1};
    outsourcing ratio \(>90\%\)) are dropped from the 2,108-firm outsourcing
    panel. This removes 233 firms, leaving \textbf{1,875 firms} in the final
    regression panel (\texttt{reg\_panel.dta}).

\end{enumerate}

\subsection*{Sample Summary}

\begin{table}[H]
  \centering
  \caption{Sample Summary Statistics}
  \label{tab:sample_summary}
  \begin{tabular}{lc}
    \toprule
    Statistic & Value \\
    \midrule
    Original sample firms                                      & 3,410 \\
    Firms with valid buy-side product codes (bianma)           & 3,376 \\
    Outsourcing firms (\texttt{invoice\_panel.dta})            & 2,108 \\
    \quad (outsourcing = same product bought \& sold)          & \\
    Regression panel firms (\texttt{reg\_panel.dta})           & 1,875 \\
    \quad (after dropping 233 intermediaries)                  & \\
    Avg.\ outsourcing products per firm (mean / median)        & 25.0 / 3 \\
    Outsourcing product types (across all firms)               & 2,143 \\
    Regression observations (T3/T6, full FE)                   & $\approx$45,669 \\
    Cities covered                                             & 262 \\
    \bottomrule
  \end{tabular}
\end{table}

%=============================================================
\section{Bridge Data Source and Matching Rate}
\label{sec:bridge}
%=============================================================

\subsection*{Purpose}

The VAT invoice database contains transaction-level information but no firm
financial data. To include \(\ln K_{ft}\) (log total assets) as a regressor in
Tables~T1, T2, T4, and T5, firm identifiers are bridged to a separate tax-filing
dataset that records balance-sheet information.

\subsection*{Bridge Path}

\[
  \underbrace{\text{Firm ID (VAT, \texttt{cid})}}_{\text{VAT database}}
  \;\longrightarrow\;
  \underbrace{\text{Tax-filing matched file}}_{\text{bridge key}}
  \;\longrightarrow\;
  \underbrace{K_{ft}\;(\text{total assets, double precision})}_{\text{tax filing data}}
\]

\subsection*{Matching Rate}

\begin{table}[H]
  \centering
  \caption{Bridge Matching Rate}
  \label{tab:match_rate}
  \begin{threeparttable}
    \begin{tabular}{lc}
      \toprule
      Item & Value \\
      \midrule
      Regression panel firms (post intermediary removal) & 1,875 \\
      Successfully matched firms (with \(\ln K\))        & $\approx$1,168 \\
      Firm-level matching rate                           & 62.3\% \\
      \midrule
      Regression observations: T1/T2/T4/T5 (require \(\ln K\)) & $\approx$29,000 \\
      Regression observations: T3/T6 (no \(\ln K\), full FE) & $\approx$45,669 \\
      Observation loss from capital requirement & $\approx$37\% \\
      \bottomrule
    \end{tabular}
  \end{threeparttable}
\end{table}

\end{document}
